\section{杂项与数学工具}
\subsection{复制线段} 
目标是在给定射线上，构造出与给定线段等长的线段。

\begin{NewMacroBox}{tkzDuplicateSegment}{\parg{pt1,pt2}\parg{pt3,pt4}\marg{pt5}}%
该宏在给定射线上构造出与给定线段等长的线段。实际上它定义了一个点。
\tkzcname{tkzDuplicateSegment} 是 \tkzcname{tkzDuplicateLen} 的新名称。

\medskip  
\begin{tabular}{lll}%
\toprule
参数             & 示例 & 说明                         \\ 

\midrule
\TAline{(pt1,pt2)(pt3,pt4)\{pt5\}} {\tkzcname{tkzDuplicateSegment}(A,B)(E,F)\{C\}}{AC=EF 且 $C \in [AB)$} \\  
\bottomrule
\end{tabular}

\medskip
\emph{宏 \tkzcname{tkzDuplicateLength} 与本宏完全相同。}
\end{NewMacroBox}

\subsubsection{\tkzcname{tkzDuplicateSegment} 的用法} 

\begin{tkzexample}[latex=6cm,small]
\begin{tikzpicture}[scale=.5]
 \tkzDefPoints{0/0/A,2/-3/B,2/5/C}
 \tkzDuplicateSegment(A,B)(A,C)  
 \tkzGetPoint{D}
 \tkzDrawSegments[new](A,B A,C)
 \tkzDrawSegment[teal](A,D)
 \tkzDrawPoints[new](A,B,C,D) 
 \tkzLabelPoints[above right=3pt](A,B,C,D)
\end{tikzpicture} 
\end{tkzexample} 

\subsubsection{用 \tkzcname{tkzDuplicateSegment} 求黄金比例} 
\begin{tkzexample}[latex=7cm,small]
\begin{tikzpicture}[rotate=-90,scale=.4]
 \tkzDefPoints{0/0/A,10/0/B}
 \tkzDefMidPoint(A,B)   
 \tkzGetPoint{I}
 \tkzDefPointWith[orthogonal,K=-.75](B,A)
 \tkzGetPoint{C}
 \tkzInterLC(B,C)(B,I)  \tkzGetSecondPoint{D}
 \tkzDuplicateSegment(B,D)(D,A) \tkzGetPoint{E}
 \tkzInterLC(A,B)(A,E)   \tkzGetPoints{N}{M}
 \tkzDrawArc[orange,delta=10](D,E)(B)
 \tkzDrawArc[orange,delta=10](A,M)(E)
 \tkzDrawLines(A,B B,C A,D)
 \tkzDrawArc[orange,delta=10](B,D)(I)
 \tkzDrawPoints(A,B,D,C,M,I)
 \tkzLabelPoints[below left](A,B,D,C,M,I)
\end{tikzpicture}
\end{tkzexample}

\subsubsection{黄金三角形或崇高三角形} 
\begin{tkzexample}[latex=7cm,small]
\begin{tikzpicture}[scale=.75] 
  \tkzDefPoints{0/0/A,5/0/C,0/5/B} 
  \tkzDefMidPoint(A,C)\tkzGetPoint{H} 
  \tkzDuplicateSegment(H,B)(H,A)\tkzGetPoint{D} 
  \tkzDuplicateSegment(A,D)(A,B)\tkzGetPoint{E} 
  \tkzDuplicateSegment(A,D)(B,A)\tkzGetPoint{G} 
  \tkzInterCC(A,C)(B,G)\tkzGetSecondPoint{F}
  \tkzDrawLine(A,C)
  \tkzDrawArc(A,C)(B) 
  \begin{scope}[arc style/.style={color=gray,%
                               style=dashed}]
    \tkzDrawArc(H,B)(D) 
    \tkzDrawArc(A,D)(B) 
    \tkzDrawArc(B,G)(F) 
  \end{scope}
  \tkzDrawSegment[dashed](H,B) 
  \tkzCompass(B,F) 
  \tkzDrawPolygon[new](A,B,F)
  \tkzDrawPoints(A,...,H)
  \tkzLabelPoints[below left](A,...,H)
\end{tikzpicture}
\end{tkzexample}

\subsection{线段长度 \tkzcname{tkzCalcLength}}
\TIKZ\ 有一个名为 \tkzname{veclen} 的函数。它
 可以计算 $AB$，其中 $A$ 和 $B$ 为两个点。

唯一的问题是，\TIKZ\ 版本在某些情况下精度不够。我的版本使用宏包 \tkzNamePack{xfp}：速度较慢，但更精确。

\begin{NewMacroBox}{tkzCalcLength}{\oarg{local options}\parg{pt1,pt2}}%
你可以用宏 \tkzcname{tkzGetLength} 保存结果；例如 \tkzcname{tkzGetLength\{dAB\}} \\
会定义宏 \tkzcname{dAB}。

\medskip
\begin{tabular}{lll}%
\toprule
参数    & 示例 & 说明       \\
\midrule
\TAline{(pt1,pt2)\{name of macro\}} {\tkzcname{tkzCalcLength}(A,B)}{\tkzcname{dAB} 以 cm 为单位给出 $AB$}
\bottomrule
\end{tabular}

\medskip
只有一个选项。

\begin{tabular}{lll}%
   
\toprule
 选项    & 默认值 & 示例       \\
\midrule
\TOline{cm}  {true}{\tkzcname{tkzCalcLength}(A,B) 在 \tkzcname{tkzGetLength\{dAB\}} 之后，\tkzcname{dAB} 以 cm 为单位给出 $AB$}
\end{tabular}
\end{NewMacroBox}

\subsubsection{圆规与正方形的构造}

\begin{tkzexample}[latex=7cm,small]
\begin{tikzpicture}[scale=1]
  \tkzDefPoint(0,0){A} \tkzDefPoint(4,0){B}
  \tkzCalcLength(A,B)\tkzGetLength{dAB}
  \tkzDefLine[perpendicular=through A](A,B)
  \tkzGetPoint{D}
  \tkzDefPointWith[orthogonal,K=-1](B,A)    
    \tkzGetPoint{F}
  \tkzGetPoint{C}
  \tkzDrawLine[add= .6 and .2](A,B)
  \tkzDrawLine(A,D) 
  \tkzShowLine[orthogonal=through A,gap=2](A,B)
  \tkzMarkRightAngle(B,A,D)
  \tkzCompasss(A,D D,C)
  \tkzDrawArc[R](B,\dAB)(80,110)
  \tkzDrawPoints(A,B,C,D)
  \tkzDrawSegments[color=gray,style=dashed](B,C C,D)
  \tkzLabelPoints[below left](A,B,C,D)
\end{tikzpicture}
\end{tkzexample}


\subsubsection{由直径确定的圆的半径}
宏 \tkzcname{tkzDefCircle[radius](A,B)} 定义一个半径，我们用 \tkzcname{tkzGetLength} 取出；结果单位为 \tkzname{cm}。

\begin{tkzexample}[latex=6cm,small]
\begin{tikzpicture}[scale=.5]
 \tkzDefPoint(0,0){A}
 \tkzDefPoint(3,-4){B}
 \tkzDefMidPoint(A,B) \tkzGetPoint{M}
 \tkzCalcLength(M,B)\tkzGetLength{rAB}
 \tkzDrawCircle(A,B)
 \tkzDrawPoints(A,B)
 \tkzLabelPoints(A,B)
 \tkzDrawSegment[dashed](A,B)
 \tkzLabelSegment(A,B){$\pgfmathprintnumber{\rAB}$}
\end{tikzpicture}
\end{tkzexample}


\subsection{从 pt 到 cm 或从 cm 到 pt 的换算}
这些宏很简单：除以或乘以 28.45274。它们是：

\begin{NewMacroBox}{tkzpttocm}{\parg{number}\marg{name of macro}}%
结果保存在一个宏中。

\medskip
\begin{tabular}{lll}%
\toprule
参数             & 示例 & 说明                         \\
\midrule
\TAline{(number)\{name of macro\}} {\tkzcname{tkzpttocm}(120)\{len\}}{\tkzcname{len} 给出以 \tkzname{cm} 为单位的数值}
\bottomrule
\end{tabular}

\medskip
你需要将 \tkzcname{len} 与单位 \tkzname{cm} 一起使用。
\end{NewMacroBox}

\subsection{单位换算} 
\begin{NewMacroBox}{tkzcmtopt}{\parg{number}\marg{name of macro}}%
结果保存在一个宏中。

\medskip
\begin{tabular}{lll}
\toprule
参数             & 示例 & 说明                         \\
\midrule
\TAline{(number)\{name of macro\}}{\tkzcname{tkzcmtopt}(5)\{len\}}{\tkzcname{len} 的长度以 \tkzname{pts} 为单位}
\bottomrule
\end{tabular}

\medskip
\emph{结果可以与 \tkzcname{len}\ \tkzname{pt} 一起使用}
\end{NewMacroBox}


\subsection{获取点的坐标}
%<--------------------------------------------------------------------------–>
%                    点的坐标 
%    结果保存在 #2x 和 #2y 中，#1 为点，我们取其坐标
% 使用 $A$ 点：\tkzGetPointCoord(A){V}，然后 \Vx = xA，\Vy = yA
% 单位 cm 
% tkzGetPointCoord 是否支持 [#1] cm 或 pt ？待办
%<--------------------------------------------------------------------------–>
\begin{NewMacroBox}{tkzGetPointCoord}{\parg{$A$}\marg{name of macro}}%
\begin{tabular}{lll}%
参数             & 示例 & 说明                         \\
\midrule
\TAline{(point)\{name of macro\}} {\tkzcname{tkzGetPointCoord}(A)\{A\}}{\tkzcname{Ax} 和 \tkzcname{Ay} 给出 $A$ 的坐标}
\end{tabular}

\medskip
\emph{将一个点的坐标保存在两个宏中。如果给定的名称为 \tkzname{p}，则 \tkzcname{px} 和 \tkzcname{py} 以 cm 为单位给出该点的坐标。}
\end{NewMacroBox}

\subsubsection{用 \tkzcname{tkzGetPointCoord} 传递坐标}

\begin{tkzexample}[width=8cm,small]
\begin{tikzpicture}
 \tkzInit[xmax=5,ymax=3]
 \tkzGrid[sub,orange]
 \tkzDrawX \tkzDrawY
 \tkzDefPoint(1,0){A}
 \tkzDefPoint(4,2){B}
 \tkzGetPointCoord(A){a}
 \tkzGetPointCoord(B){b}
 \tkzDefPoint(\ax,\ay){C}
 \tkzDefPoint(\bx,\by){D}
 \tkzDrawPoints[color=red](C,D)
\end{tikzpicture}
\end{tkzexample}

\subsubsection{用 \tkzcname{tkzGetPointCoord} 求向量之和}
\begin{tkzexample}[width=6cm,small]
\begin{tikzpicture}[>=latex]
  \tkzDefPoint(1,4){a}
  \tkzDefPoint(3,2){b}
  \tkzDefPoint(1,1){c}
  \tkzDrawSegment[->,red](a,b)
  \tkzGetPointCoord(c){c}
  \draw[color=blue,->](a) -- ([shift=(b)]\cx,\cy) ;
  \draw[color=purple,->](b) -- ([shift=(b)]\cx,\cy) ;
  \tkzDrawSegment[->,blue](a,c)
  \tkzDrawSegment[->,purple](b,c)
\end{tikzpicture}
\end{tkzexample}

\subsection{交换点的标签}

\begin{NewMacroBox}{tkzSwapPoints}{\parg{$pt1$,$pt2$}}%
\begin{tabular}{lll}%
参数             & 示例 & 说明                         \\
\midrule
\TAline{(pt1,pt2)} {\tkzcname{tkzSwapPoints}(A,B)}{现在 $A$ 具有 $B$ 的坐标}
\end{tabular}

\emph{两点交换了各自的坐标。}
\end{NewMacroBox}

\subsubsection{\tkzcname{tkzSwapPoints} 的用法}

\begin{tkzexample}[width=6cm,small]
\begin{tikzpicture}
  \tkzDefPoints{0/0/O,5/-1/A,2/2/B}
   \tkzSwapPoints(A,B)
   \tkzDrawPoints(O,A,B)
   \tkzLabelPoints(O,A,B)
\end{tikzpicture}
\end{tkzexample}

\subsection{点积}
在欧几里得几何中，由笛卡尔坐标计算的两个向量的点积被广泛使用。

\begin{NewMacroBox}{tkzDotProduct}{\parg{$pt1$,$pt2$,$pt3$}}%
  两个向量 $\overrightarrow{u} = [a,b]$ 和  $\overrightarrow{v} = [a',b']$ 的点积定义为：$\overrightarrow{u}\cdot \overrightarrow{v} = aa' + bb'$

$\overrightarrow{u} = \overrightarrow{pt1pt2}$ $\overrightarrow{v} = \overrightarrow{pt1pt3}$
  
\begin{tabular}{lll}%
参数             & 示例 & 说明                         \\
\midrule
\TAline{(pt1,pt2,pt3)} {\tkzcname{tkzDotProduct}(A,B,C)}{结果为 $\overrightarrow{AB}\cdot \overrightarrow{AC}$}
\end{tabular}

\emph{结果是一个数值，可以用 \tkzcname{tkzGetResult} 取出。}
\end{NewMacroBox}

\subsubsection{简单示例} % (fold)
\label{ssub:simple_example}

\begin{tkzexample}[small,latex=7cm]
\begin{tikzpicture}
  \tkzDefPoints{-2/-3/A,4/0/B,1/3/C}
  \tkzDefPointBy[projection= onto A--B](C)  
  \tkzGetPoint{H}
  \tkzDrawSegment(C,H)
  \tkzMarkRightAngle(C,H,A)
  \tkzDrawSegments[vector style](A,B A,C)
  \tkzDrawPoints(A,H) \tkzLabelPoints(A,B,H)
  \tkzLabelPoints[above](C)
  \tkzDotProduct(A,B,C) \tkzGetResult{pabc}
 % \pgfmathparse{round(10*\pabc)/10}
  \let\pabc\pgfmathresult
  \node at (1,-3) {$\overrightarrow{PA}\cdot \overrightarrow{PB}=\pabc$};
  \tkzDotProduct(A,H,B) \tkzGetResult{phab}
 % \pgfmathparse{round(10*\phab)/10}
  \let\phab\pgfmathresult
  \node at (1,-4) {$PA \times PH = \phab $};
\end{tikzpicture}
\end{tkzexample}
% subsubsection simple_example (end)


\subsubsection{共圆点} % (fold)
\label{ssub:cocyclicpts}

\begin{tkzexample}[small,latex=7cm]
\begin{tikzpicture}[scale=.75]
  \tkzDefPoints{1/2/O,5/2/B,2/2/P,3/3/Q}
  \tkzInterLC[common=B](O,B)(O,B) \tkzGetFirstPoint{A}
  \tkzInterLC[common=B](P,Q)(O,B) \tkzGetPoints{C}{D}
  \tkzDrawCircle(O,B)
  \tkzDrawSegments(A,B C,D)
  \tkzDrawPoints(A,B,C,D,P)
  \tkzLabelPoints(P)
  \tkzLabelPoints[below left](A,C)
  \tkzLabelPoints[above right](B,D)
  \tkzDotProduct(P,A,B) \tkzGetResult{pab}
  \pgfmathparse{round(10*\pab)/10}
  \let\pab\pgfmathresult
  \tkzDotProduct(P,C,D) \tkzGetResult{pcd}
  \pgfmathparse{round(10*\pcd)/10}
  \let\pcd\pgfmathresult
  \node at (1,-3) {%
  $\overrightarrow{PA}\cdot \overrightarrow{PB} =
   \overrightarrow{PC}\cdot \overrightarrow{PD}$};
  \node at (1,-4)%
  {$\overrightarrow{PA}\cdot \overrightarrow{PB}=\pab$};
 \node at (1,-5){%
 $\overrightarrow{PC}\cdot \overrightarrow{PD} =\pcd$};
\end{tikzpicture}
\end{tkzexample}
% subsubsection cocyclicpts (end)

\newpage
\subsection{点关于圆的幂}

\begin{NewMacroBox}{tkzPowerCircle}{\parg{$pt1$}\parg{$pt2$,$pt3$}}%
\begin{tabular}{lll}%
参数             & 示例 & 说明                         \\
\midrule
\TAline{(pt1)(pt2,pt3)} {\tkzcname{tkzPowerCircle}(A)(O,M)}{$A$ 关于圆 (O,A) 的幂}
\end{tabular}

\emph{结果是一个数值，表示一个点关于一个圆的幂。}
\end{NewMacroBox}

\subsubsection{由根轴求幂} % (fold)
\label{ssub:power}

在本例中，已画出根轴 $(EF)$。在 $(EF)$ 上选取一点 $H$，并计算了 $H$ 关于圆心为 $A$ 的圆的幂，以及 $HS^2$。你可以验证，$H$ 关于圆心为 $C$ 的圆的幂，以及 $HS'^2$、$HT^2$、$HT'^2$ 给出相同的结果。

\begin{tkzexample}[small,latex=7cm]
\begin{tikzpicture}[scale=.5]
  \tkzDefPoints{-1/0/A,0/5/B,5/-1/C,7/1/D}
  \tkzDrawCircles(A,B C,D)
  \tkzDefRadicalAxis(A,B)(C,D) \tkzGetPoints{E}{F}
  \tkzDrawLine[add=1 and 2](E,F)
  \tkzDefPointOnLine[pos=1.5](E,F) 
  \tkzGetPoint{H}
  \tkzDefLine[tangent from = H](A,B)
  \tkzGetPoints{T}{T'}
  \tkzDefLine[tangent from = H](C,D)
  \tkzGetPoints{S}{S'}
  \tkzDrawSegments(H,T H,T' H,S H,S')
  \tkzDrawPoints(A,B,C,D,E,F,H,T,T',S,S')
  \tkzPowerCircle(H)(A,B) \tkzGetResult{pw}
  \tkzDotProduct(H,S,S) \tkzGetResult{phtt}
  \node {Power $\approx \pw \approx \phtt$};
\end{tikzpicture}
\end{tkzexample}
% subsubsection power (end)

\subsection{根轴}

在几何中，两个不同心圆的根轴是关于两个圆具有相同幂的点的集合。这里 |\tkzDefRadicalAxis(A,B)(C,D)| 给出两个圆 $\mathcal{C}(A,B)$ 和 $\mathcal{C}(C,D)$ 的根轴。 

\begin{NewMacroBox}{tkzDefRadicalAxis}{\parg{$pt1$,$pt2$}\parg{$pt3$,$pt4$}}%
\begin{tabular}{lll}%
参数             & 示例 & 说明                         \\
\midrule
\TAline{(pt1,pt2)(pt3,pt4)} {\tkzcname{tkzDefRadicalAxis}(A,B)(C,D)}{圆心为 $A$ 和 $C$ 的两个圆}
\midrule
\end{tabular}


\emph{结果为根轴上的两个点。}
\end{NewMacroBox}

\subsubsection{两个相离的圆} % (fold)
\label{ssub:two_circles_disjointed}


\begin{tkzexample}[small,latex=8cm]
\begin{tikzpicture}[scale=.75]
    \tkzDefPoints{-1/0/A,0/2/B,4/-1/C,4/0/D}
    \tkzDrawCircles(A,B C,D)
    \tkzDefRadicalAxis(A,B)(C,D)
     \tkzGetPoints{E}{F}
    \tkzDrawLine[add=1 and 2](E,F)
    \tkzDrawLine[add=.5 and .5](A,C)
\end{tikzpicture}
\end{tkzexample}
% subsubsection two_circles_disjointed (end)

\subsubsection{两个相交的圆} % (fold)
\label{sub:two_intersecting_circles}
\begin{tkzexample}[small,latex=8cm]
\begin{tikzpicture}[scale=.5]
  \tkzDefPoints{-1/0/A,0/2/B,3/-1/C,3/-2/D}
  \tkzDrawCircles(A,C B,D)
  \tkzDefRadicalAxis(A,C)(B,D)
  \tkzGetPoints{E}{F}
  \tkzDrawPoints(A,B,C,D,E,F)
    \tkzLabelPoints(A,B,C,D,E,F)
  \tkzDrawLine[add=.25 and .5](E,F)
  \tkzDrawLine[add=.25 and .25](A,B)
\end{tikzpicture}
\end{tkzexample}
% subsection two_intersecting_circles (end)


\subsubsection{两个外切圆} % (fold)
\label{sub:two_externally_tangent_circles}
\begin{tkzexample}[small,latex=8cm]
\begin{tikzpicture}[scale=.5]
  \tkzDefPoints{0/0/A,4/0/B,6/0/C}
  \tkzDrawCircles(A,B C,B)
  \tkzDefRadicalAxis(A,B)(C,B)
  \tkzGetPoints{E}{F}
  \tkzDrawPoints(A,B,C,E,F)
    \tkzLabelPoints(A,B,C,E,F)
  \tkzDrawLine[add=1 and 1](E,F)
  \tkzDrawLine[add=.5 and .5](A,B)
\end{tikzpicture}
\end{tkzexample}
% subsection two_externally_tangent_circles (end)


\subsubsection{两个内切圆} % (fold)
\label{sub:deux_cercles_tangents_interieurement}
\begin{tkzexample}[small,latex=8cm]
\begin{tikzpicture}[scale=.5]
  \tkzDefPoints{0/0/A,3/0/B,5/0/C}
  \tkzDrawCircles(A,C B,C)
  \tkzDefRadicalAxis(A,C)(B,C)
  \tkzGetPoints{E}{F}
  \tkzDrawPoints(A,B,C,E,F)
  \tkzLabelPoints[below right](A,B,C,E,F)
  \tkzDrawLine[add=1 and 1](E,F)
  \tkzDrawLine[add=.5 and .5](A,B)
\end{tikzpicture}
\end{tkzexample}
% subsection deux_cercles_tangents_interieurement (end)

\subsubsection{三个圆} % (fold)
\label{ssub:threecircles}



\begin{tkzexample}[small,latex=7cm]
\begin{tikzpicture}[scale=.4]
  \tkzDefPoints{0/0/A,5/0/a,7/-1/B,3/-1/b,5/-4/C,2/-4/c}
  \tkzDrawCircles(A,a B,b C,c)
  \tkzDefRadicalAxis(A,a)(B,b) \tkzGetPoints{i}{j}
  \tkzDefRadicalAxis(A,a)(C,c) \tkzGetPoints{k}{l}
  \tkzDefRadicalAxis(C,c)(B,b) \tkzGetPoints{m}{n}
  \tkzDrawLines[new](i,j k,l m,n)
\end{tikzpicture}
\end{tkzexample}
% subsubsection threecircles (end)
  
\subsection{\tkzcname{tkzIsLinear}、\tkzcname{tkzIsOrtho}}
 \begin{NewMacroBox}{tkzIsLinear}{\parg{$pt1$,$pt2$,$pt3$}}%
 \begin{tabular}{lll}%
 参数             & 示例 & 说明                         \\
 \midrule
 \TAline{(pt1,pt2,pt3)} {\tkzcname{tkzIsLinear}(A,B,C)}{$A,B,C$ 共线？}
 \midrule
 \end{tabular}
 
 \emph{\tkzcname{tkzIsLinear} 测试三个点 $pt1$、$pt2$、$pt3$ 是否共线。测试允许约 $0.06^\circ$ 的角度误差。}
 \end{NewMacroBox}
 
 \begin{NewMacroBox}{tkzIsOrtho}{\parg{$pt1$,$pt2$,$pt3$}}%
 \begin{tabular}{lll}%
 参数             & 示例 & 说明                         \\
 \midrule
 \TAline{(pt1,pt2,pt3)} {\tkzcname{tkzIsOrtho}(A,B,C)}{$(AB)\perp (AC)$ ？}
 \midrule
 \end{tabular}
 
 \emph{\tkzcname{tkzIsOrtho} 测试直线 $(pt1pt2)$ 与 $(pt1pt3)$ 是否正交，容差相同。}
 \end{NewMacroBox}
 
 \subsubsection{\tkzcname{tkzIsOrtho} 与 \tkzcname{tkzIsLinear} 的用法}
   
\begin{tkzexample}[small,latex=7cm]
  \begin{tikzpicture}
  \tkzDefPoints{1/-2/A,5/0/B}
  \tkzDefCircle[diameter](A,B) \tkzGetPoint{O}
  \tkzDrawCircle(O,A)
  \tkzDefPointBy[rotation= center O angle 60](B) 
  \tkzGetPoint{C}
  \tkzDefPointBy[rotation= center O angle 60](A) 
  \tkzGetPoint{D}
  \tkzDrawCircle(O,A)
  \tkzDrawPoints(A,B,C,D,O)
  \tkzIsOrtho(C,A,B)
  \iftkzOrtho
    \tkzDrawPolygon[blue](A,B,C)
  \tkzDrawPoints[blue](A,B,C,D)
  \else
  \tkzDrawPoints[red](A,B,C,D)
  \fi
   \tkzIsLinear(O,C,D)
    \iftkzLinear
    \tkzDrawSegment[orange](C,D)
    \fi
\end{tikzpicture}
  
\end{tkzexample}

  
\endinput
